\begin{textsample}{POS Dim 1 – pe – Score 42.00 – pe/pe\_ef\_22.txt}  \label{ex:f1_pos_244}
3.3.4 EIXOS ESTRUTURANTES : AS PRÁTICAS DE LINGUAGEM E OS CAMPOS DE ATUAÇÃO Assim como na Base Nacional Comum Curricular/BNCC e em documentos curriculares estaduais anteriores ( BCC, OTM, PCE³ ), os objetos de conhecimento e as habilidades, neste currículo, estruturam -se dentro das práticas de linguagem : leitura/escuta, produção escrita, oralidade e, agora, \textbf{multissemiótica} e \textbf{análise} linguística/semiótica ( que envolve os conhecimentos sobre o sistema alfabético de escrita, a norma padrão, os aspectos textuais, \textbf{discursivos}, gramaticais, os modos de organização e os elementos de outras semioses, entre outros ).
Desse modo, > [... ] a língua somente poderá ser entendida como uma ação contextualizada e historicamente situada ; sempre inserida numa situação particular de interação e, portanto, nunca inteiramente despregada das condições concretas de uma determinada prática social, não podendo, assim, ser avaliada senão em situação. ( BCC, 2008, p. 67 ).
Dentro da concepção de língua aqui defendida, as práticas de linguagem/os eixos de ensino se \textbf{materializam} em práticas historicamente situadas/contextualizadas, isto é, em campos/esferas \textbf{discursivas} ; relacionam -se às práticas de uso e reflexão, e são \textbf{influenciadas} pelas condições de produção e \textbf{recepção} dos textos.
No eixo da Leitura, o \textbf{foco} recai sobre a interação ativa entre leitor/ouvinte/espectador e os textos, \textbf{tanto} na modalidade escrita quanto oral, além dos recursos semióticos de diferentes esferas \textbf{discursivas}, com a finalidade de proporcionar o contato dos estudantes com diferentes experiências leitoras e estratégias de leitura para torná -lo um leitor proficiente e \textbf{crítico}.
Trata -se, portanto, de uma prática social que \textbf{implica}, necessariamente, a \textbf{compreensão} na qual a construção de sentido vai se constituindo antes da leitura propriamente \textbf{dita} ( BRASIL, 1997 ), e pelo estabelecimento do profícuo diálogo e da negociação entre os interlocutores. ³ As siglas referem -se, respectivamente, a Base Curricular Comum para as Redes Públicas de Ensino de PE, Orientações Teórico-Metodológicas e \textbf{Parâmetros} Curriculares para a Educação Básica do Estado de PE.
Nessa perspectiva, ensinar a ler é mostrar aos estudantes que é preciso considerar os contextos de produção ( inclusive a literária ) em que as interações sociais acontecem, bem como reconhecer a importância das culturas do escrito e interpretar imagens e recursos semióticos que constituem muitos gêneros \textbf{digitais}.
É concorrer para o desenvolvimento da \textbf{capacidade} de relacionar textos e diferentes linguagens, além de permitir a interação com variadas crenças, valores, concepções, conflitos, subjetividades e identidades, possibilitando o autoconhecimento e o desenvolvimento de uma postura respeitosa diante daquilo que é diferente, entre outros aspectos.
Para \textbf{tanto}, é \textbf{imprescindível} disponibilizar para os estudantes, ao longo da educação básica, o contato sistemático com uma variedade de gêneros, exemplares de textos, suportes, procedimentos de leitura, contextos de produção, visando, essencialmente, a ( 1 ) reconstruir e avaliar as condições de produção e \textbf{recepção} dos textos de diferentes gêneros e esferas \textbf{discursivas}, como também a textualidade – organização, \textbf{progressão} temática, estabelecimento de relações entre as partes do texto – ; ( 2 ) avaliar criticamente as temáticas tratadas e a validade das informações e dados ; ( 3 ) compreender os efeitos de sentido provocados pelos usos de recursos linguísticos e \textbf{multissemióticos} ; ( 4 ) selecionar procedimentos de leitura adequados a diferentes objetivos e interesses e ( 5 ) envolver -se pela leitura de textos literários, textos de divulgação científica e textos jornalísticos.
O eixo da produção textual compreende as práticas de linguagem relativas à interação por meio de textos escritos, orais e \textbf{multissemióticos} com diferentes propósitos comunicativos e nos diversos campos de atividade humana, em função dos quais o sentido para o que se escreve é construído.
Produzir linguagens significa produzir discursos.
É \textbf{dizer} algo para alguém, de uma determinada forma, num contexto.
As escolhas feitas ao produzir um discurso não são aleatórias, mas decorrentes das condições em que esse discurso é realizado ( BRASIL, 1997 ).
Dessa forma, a produção e \textbf{recepção} de textos passam a \textbf{preocupar} -se com os contextos ( finalidades, interlocutores, suportes, recursos ) e com os processos de produção ( planejamento, textualização e revisão, reescrita/edição ), corroborando com a concepção de língua enquanto sistema forjado histórica e socialmente pela ação dos usuários e de produção de texto enquanto processo/lugar de interação e não \textbf{mero} produto.
As estratégias didáticas com esse eixo devem iniciar sempre com propostas de produção de textos ( orais e escritos ) significativas e contextualizadas.
Com base em Beaugrande, Marcuschi ( 2008 ) defende que a situacionalidade é um critério da textualidade importante, \textbf{tanto} para interpretação como para orientação da produção textual ; inclusive, deve ser considerada para avaliação do texto : revisão, reescrita e \textbf{edição}.
Isto \textbf{posto}, é preciso oportunizar os estudantes a vivenciar esse jogo \textbf{complexo} de produção de sentido e assumir -se como autor/coautor dos discursos, sendo capaz, sobretudo de ( 1 ) analisar os diferentes contextos e situações sociais nos quais os textos são produzidos e sobre as diferenças – em termos formais, estilísticos e linguísticos – que tais contextos determinam ; ( 2 ) relacionar diferentes textos ; ( 3 ) selecionar os procedimentos de escrita adequados aos contextos de produção, além das informações, dados, \textbf{argumentos} e outras referências a serviço do projeto textual ; ( 4 ) textualizar, ou seja, organizar as informações e \textbf{progressão} temática, estabelecer relações entre as partes do texto, usar – a serviço do projeto textual e da produção de sentido – os recursos linguísticos e \textbf{multissemióticos} mais apropriados e ( 5 ) empregar, de acordo com o contexto de interação, os aspectos notacionais e gramaticais.
O eixo da Oralidade compreende as práticas de linguagem na modalidade falada da língua em situações de uso oral da linguagem.
Nesse sentido, ensinar a oralidade \textbf{implica} conceber a língua oral como objeto autônomo de ensino ( SCHNEUWLY ; DOLZ, 2010 ), considerando as dimensões \textbf{discursivas} e materiais \textbf{pertinentes} a essa modalidade, com destaque para o “ como \textbf{dizer} e o que se tem a \textbf{dizer} ”, bem como refletir linguisticamente em torno dessa modalidade em diversos contextos de interação, principalmente, em situações mais formais do que aquelas já vivenciadas pelos estudantes.
O ensino e a aprendizagem da oralidade como prática social visam a desenvolver \textbf{capacidades} de ação linguístico-discursivas e a ampliar o repertório e a visão de língua dos falantes, o domínio da variedade padrão, a \textbf{compreensão} das relações entre fala e escrita, o sucesso escolar dos estudantes, a superação de preconceitos, entre outros.
Logo, o estudo das características e finalidade dos diferentes gêneros textuais orais e seus modos de realização será objeto \textbf{central} desse eixo, sem \textbf{perder} de \textbf{vista} a oralização de textos escritos, que, diferentemente da oralidade, trata -se de uma forma de atividade com o apoio do texto escrito.
A realização do trabalho pedagógico com a oralidade na escola e a formação de estudantes competentes linguisticamente também na modalidade oral da língua compreende ( 1 ) analisar os diferentes contextos e situações sociais em que se produzem textos orais e as diferenças em termos formais, estilísticos e linguísticos que tais contextos determinam ; ( 2 ) compreender textos orais ; ( 3 ) produzir textos orais pertencentes a gêneros orais diversos e ( 4 ) compreender os efeitos de sentido provocados pelos usos de recursos linguísticos e \textbf{multissemióticos} e, especialmente, ( 4 ) estabelecer um ambiente favorável à manifestação dos pensamentos, dos sentimentos e das identidades, acolhendo as diferenças.
Por fim, o eixo da Análise linguística/semiótica, assim como já consolidado em documentos anteriores, vincula -se à perspectiva do uso-reflexão-uso⁴ da língua e a serviço das práticas propostas nos eixos de Oralidade, Leitura e Produção de textos.
A finalidade é que os estudantes reflitam sobre as diferentes possibilidades e recursos da língua na produção de sentido e adequado ao contexto de interação.
Esse eixo contempla a \textbf{análise} e reflexão sobre os textos orais, escritos e \textbf{multissemióticos}, envolvendo os conhecimentos sobre o sistema alfabético de escrita, a norma padrão, os aspectos \textbf{discursivos}, textuais, gramaticais, os modos de organização linguística ( fonética, fonológica, conhecimentos grafônicos, ortográficos e lexicais, morfossintática, semântica e pragmática ), além dos elementos de outras semioses.
Sempre no \textbf{intuito} da produção e \textbf{compreensão} de sentidos \textbf{materializados} nos mais diversos gêneros textuais.
Com relação à composição dos textos multissemióticos/multimodais, à semelhança da BNCC ( 2017 ), este documento também orienta quanto ao estilo e formas de organização das linguagens que os integram.
O eixo de reflexão sobre a língua e outras semioses contempla \textbf{análise} e avaliação sobre as \textbf{inter-relações} provenientes entre : a ) as imagens visuais estáticas ( cor, intensidade, \textbf{foco}, \textbf{profundidade}, plano/ângulo/lado ), b ) as imagens dinâmicas e performances ( movimento, duração, sincronização, montagem, ritmo, tipo de movimento ) e c ) música ( instrumentos, timbres, melodia, harmonia, ritmo ).
Esse tratamento das multimodalidades e multissemioses em textos que integram o mundo \textbf{contemporâneo} \textbf{exigem} práticas de \textbf{compreensão} e produção que corrobora para o \textbf{multiletramento}.
Na \textbf{inter-relação} com os demais eixos, o trabalho com a \textbf{análise} linguística/semiótica, priorizando os conhecimentos significativos para o cotidiano dos usuários da língua, precisa ser organizado partindo de uma \textbf{abordagem} epilinguística para a metalinguística.
Compreendendo a epilinguística como “ a prática que \textbf{opera} sobre a própria linguagem, compara as expressões, transforma -as, experimenta novos modos de construção canônicos ou não, brinca com a linguagem, \textbf{investe} as formas linguísticas de novas significações ” ( FRANCHI, 2006 apud BAGNO, 2010, p. 94 ) ; enquanto a metalinguística ( do grego metá : do outro \textbf{lado} de, para além de ) são aplicadas “ para designar o conjunto de recursos fonomorfossintáticos e lexicais que, acionados pela semântica e pela pragmática, usamos para nossa interação sociocomunicativa com os outros seres humanos no mundo ” ( BAGNO, 2010, p. 88 ).
Essa \textbf{abordagem} visa favorecer uma pedagogia de reflexão sobre as várias possibilidades de operacionalizar os recursos expressivos presentes nos textos escritos e orais ( incluindo os multimodais ).
Por entender que os eixos apresentados se \textbf{materializam} em práticas de linguagem situadas, tais eixos se articulam com a \textbf{categoria} campos de atuação, \textbf{sinalizando}, conforme \textbf{dito} anteriormente, para a necessidade de contextualizar os estudos sobre a língua, isto é, essas práticas são próprias de situações da vida em sociedade, devendo ser postas, portanto, em contextos significativos para os estudantes.
Cabe destacar que os campos de atuação orientam a seleção de gêneros, práticas e procedimentos em cada uma das esferas \textbf{discursivas}, além de assumirem a função de organizar os conhecimentos sobre a língua e outras linguagens ( função didática ), considerando que as fronteiras entre os campos mostram -se tênues, permitindo o trânsito de gêneros textuais e a intersecção de diferentes \textbf{maneiras}.
Para \textbf{tanto}, tal qual a BNCC, cinco campos foram selecionados – campo da vida cotidiana, campo artístico-literário, campo das práticas de estudo e pesquisa, campo jornalístico-midiático e campo de atuação na vida pública – considerando a importância deles para o uso da linguagem dentro e fora da escola, para o exercício da cidadania, para o desenvolvimento da autonomia e do pensamento \textbf{crítico}, para elaboração/consolidação de um projeto de vida e sociedade, para apreciação/fruição estética e simbólica da linguagem, para o exercício da empatia e do diálogo, entre outras dimensões formativas não menos importantes.
No \textbf{quadro} abaixo, são apresentados a distribuição dos campos nas etapas do ensino fundamental e uma breve caracterização dos mesmos. ⁴ Perspectiva apresentada nos \textbf{Parâmetros} Curriculares Nacionais de Língua Portuguesa os quais conceberam que \textbf{tanto} o ponto de partida quanto a finalidade do ensino da língua é a produção/compreensão dos discursos ( usos ).
Assim, as atividades didáticas precisariam ser organizadas em função da \textbf{análise} das produções ( seja no papel de leitor e/ou produtor de textos ) e do próprio processo de produção.

% matched lemmas: abordagem, análise, argumento, capacidade, categoria, central, complexo, compreensão, contemporâneo, crítico, digital, discursivo, dizer, edição, exigir, foco, implicar, imprescindível, influenciar, inter-relação, intuito, investir, lado, maneira, materializar, mero, multiletramento, multissemiótico, operar, parâmetro, perder, pertinente, postar, preocupar, profundidade, progressão, quadro, recepção, sinalizar, tanto, vista
\end{textsample}
